% !PRINTED 350
\weberpage{350}
%<<CONT

%% [running head, not body text -- commented out 2026-08-07] Page furniture (running
%% head, folio, signature) is excluded from this edition by standing policy.
%% \noindent 350\hfill Elfter Abschnitt.\hfill \S\ 95.

Bei positiver Diskriminante gibt es in jeder Formklasse Formen mit positiven ersten Koeffizienten, und daher entsprechen sich Formenklassen und Idealklassen in eindeutiger Weise.

Eine Basis des Hauptideals $\eta$ ist $(\eta, \eta\theta)$ und
\[
\left.\begin{aligned}
\theta &= \frac{\sqrt{\Delta}}{2}, & \theta^2 &= \frac{\Delta}{4}, & \Delta &\equiv 0\\
\theta &= \frac{-1+\sqrt{\Delta}}{2}, & \theta^2 &= \frac{\Delta-1}{4} - \theta, & \Delta &\equiv 1
\end{aligned}\right\}\ (\mathrm{mod}\ 4).
\]
Sind $x, y$ ganze rationale Zahlen und
\[
\eta = x + y\theta,
\]
so ist diese Basis in den beiden Fällen:
\[
\begin{aligned}
(\alpha_1, \alpha_2) &= \left(x + y\theta,\ \frac{\Delta y}{4} + \theta x\right),\\
&= \left(x + y\theta,\ y\,\frac{\Delta-1}{4} + (x-y)\theta\right)
\end{aligned}
\]
und die Substitutionsdeterminante $A$ [\S\ 87, (1)] ist
\[
x^2 - \frac{\Delta}{4}y^2, \quad x^2 - xy + \frac{1-\Delta}{4}y^2 = N(\eta).
\]
Die Basis $(\alpha_1, \alpha_2)$ ist also positiv, wenn die Norm von $\eta$ positiv ist (sonst müßte $\alpha_1$ mit $\alpha_2$ vertauscht werden). Aus der Basisform
\[
\lambda = \alpha_1 t_1 + \alpha_2 t_2 = \eta(t_1 + \theta t_2)
\]
erhält man
\[
N(\lambda) = N(\eta)\,T,
\]
und danach wird $T$ die \gesperrt{Hauptform}
\weq{9}{
(1,\,0,\,\tfrac{\Delta}{4}), \quad \left(1,\,-1,\,\frac{+1-\Delta}{4}\right).
}
Wäre $N(\eta)$ negativ, so würde man für $T$ die Formen
\[
\left(-1,\,0,\,\frac{-\Delta}{4}\right), \quad \left(-1,\,+1\ \frac{-1-\Delta}{4}\right) % [sic] -- no comma printed before the fraction
\]
erhalten haben, und diese sind mit (9) nur dann äquivalent, wenn die \gesperrt{Pell}sche Gleichung $t^2 - \Delta u^2 = -4$ lösbar ist.

\rule{4cm}{0.4pt}
